% GENERATED FILE. Edit canonical lessons, then run npm run generate:books.
% canonical-source-hash: fnv1a64:71938a296848a03c
% canonical-lessons: MR-C527-cirne, MR-C527-ukalne-boil, MR-C527-talne-fry, MR-C527-dhavalne, MR-C527-karapne, MR-R527-review-520-527

\chapter{At the Stove}
\label{ch:mr-at-the-stove}
\hlchaptermodality{527}

\begin{chapteropening}
\textbf{By the end of this chapter:} \emph{I can say how food is chopped, boiled, fried and stirred.}

Complete the last lesson of chapter 527: I can say how food is chopped, boiled, fried and stirred.
\end{chapteropening}

\section[\texorpdfstring{\mr{चिरणे} (cirṇe)}{चिरणे (cirṇe)}]{\mr{चिरणे} (cirṇe) — to chop; to slice}

\label{lesson:MR-C527-cirne}



\begin{quote}
Before the new one: say the Marathi for to knead, then the Marathi for to grate.
\end{quote}

\subsection*{You'll want to know: \mr{चिरणे}}
\textbf{\mr{चिरणे}} — \emph{cirṇe} — \textquotedblleft{}to chop; to slice\textquotedblright{}.

\begin{quote}
\mr{भाजीसाठी} \mr{मी} \mr{दोन} \mr{कांदे} \mr{चिरले}.
\end{quote}

\begin{quote}
\emph{bhājīsāṭhī mī don kānde cirle.}
\end{quote}

\begin{quote}
I chopped two onions for the vegetable dish.
\end{quote}

Past, with -\mr{ने} (he …): \mr{त्याने} \mr{कांदा} \mr{चिरला}, \emph{tyāne kāndā cirlā}.

\begin{cousinweb}
\textbf{\mr{चिरणे}} (\emph{cirṇe}) is inherited: CDIAL lists Marathi \emph{cirṇẽ}, \textquotedblleft{}to split\textquotedblright{}, under \emph{*cīrayati}, \textquotedblleft{}splits\textquotedblright{} (CDIAL 4844).
\end{cousinweb}

\begin{grammarlens}[title={What you've built}]
One more word for describing what you lived and felt.
\end{grammarlens}

\subsection*{Your turn}
\begin{itemize}
\raggedright
  \item \emph{Say it:} \emph{cirṇe}
  \item \emph{Say it:} \emph{cirṇe}, once more
  \item \emph{Say it:} \emph{bhājīsāṭhī mī don kānde cirle.} — I chopped two onions for the vegetable dish.
  \item \emph{Recall:} say the Marathi for to knead, then the Marathi for to grate, then say \emph{cirṇe} again
\end{itemize}

\begin{tcolorbox}[breakable,colback=teal!4,colframe=teal!35!black,title={Before you move on}]
What does \mr{चिरणे} mean? (\textquotedblleft{}To chop; to slice\textquotedblright{}.) Say it once more.
\end{tcolorbox}
\section[\texorpdfstring{\mr{उकळणे} (ukaḷṇe)}{उकळणे (ukaḷṇe)}]{\mr{उकळणे} (ukaḷṇe) — to boil}

\label{lesson:MR-C527-ukalne-boil}



\begin{quote}
Before the new one: say the Marathi for to grate, then the Marathi for to chop.
\end{quote}

\subsection*{You'll want to know: \mr{उकळणे}}
\textbf{\mr{उकळणे}} — \emph{ukaḷṇe} — \textquotedblleft{}to boil\textquotedblright{}.

\begin{quote}
\mr{चहासाठी} \mr{आईने} \mr{पाणी} \mr{उकळलं}.
\end{quote}

\begin{quote}
\emph{cahāsāṭhī āīne pāṇī ukaḷlaṁ.}
\end{quote}

\begin{quote}
Mother boiled water for tea.
\end{quote}

Past, with -\mr{ने} (he …): \mr{त्याने} \mr{पाणी} \mr{उकळलं}, \emph{tyāne pāṇī ukaḷlaṁ}.

Both ways: \mr{पाणी} \mr{उकळतं}, the water boils; \mr{मी} \mr{पाणी} \mr{उकळतो}, I boil the water.

\begin{cousinweb}
\textbf{\mr{उकळणे}} (\emph{ukaḷṇe}) is inherited: CDIAL lists Marathi \emph{ukaḷṇẽ}, \textquotedblleft{}to boil\textquotedblright{}, under \emph{*utkalati}, \textquotedblleft{}goes out or up\textquotedblright{} (CDIAL 1716).
\end{cousinweb}

\begin{grammarlens}[title={What you've built}]
One more word for describing what you lived and felt.
\end{grammarlens}

\subsection*{Your turn}
\begin{itemize}
\raggedright
  \item \emph{Say it:} \emph{ukaḷṇe}
  \item \emph{Say it:} \emph{ukaḷṇe}, once more
  \item \emph{Say it:} \emph{cahāsāṭhī āīne pāṇī ukaḷlaṁ.} — Mother boiled water for tea.
  \item \emph{Recall:} say the Marathi for to grate, then the Marathi for to chop, then say \emph{ukaḷṇe} again
\end{itemize}

\begin{tcolorbox}[breakable,colback=teal!4,colframe=teal!35!black,title={Before you move on}]
What does \mr{उकळणे} mean? (\textquotedblleft{}To boil\textquotedblright{}.) Say it once more.
\end{tcolorbox}
\section[\texorpdfstring{\mr{तळणे} (taḷṇe)}{तळणे (taḷṇe)}]{\mr{तळणे} (taḷṇe) — to fry}

\label{lesson:MR-C527-talne-fry}



\begin{quote}
Before the new one: say the Marathi for to chop, then the Marathi for to boil.
\end{quote}

\subsection*{You'll want to know: \mr{तळणे}}
\textbf{\mr{तळणे}} — \emph{taḷṇe} — \textquotedblleft{}to fry\textquotedblright{}.

\begin{quote}
\mr{पाहुण्यांसाठी} \mr{आजीने} \mr{गरम} \mr{पापड} \mr{तळले}.
\end{quote}

\begin{quote}
\emph{pāhuṇyāṁsāṭhī ājīne garam pāpaḍ taḷle.}
\end{quote}

\begin{quote}
Grandmother fried hot papads for the guests.
\end{quote}

Past, with -\mr{ने} (he …): \mr{त्याने} \mr{पापड} \mr{तळला}, \emph{tyāne pāpaḍ taḷlā}.

\begin{cousinweb}
\textbf{\mr{तळणे}} (\emph{taḷṇe}) is inherited: CDIAL lists Marathi \emph{taḷṇẽ}, \textquotedblleft{}to fry, cook in ghee\textquotedblright{}, under \emph{*talati}, \textquotedblleft{}fries\textquotedblright{} (CDIAL 5736), beside Hindi \emph{talnā}.
\end{cousinweb}

\begin{grammarlens}[title={What you've built}]
One more word for describing what you lived and felt.
\end{grammarlens}

\subsection*{Your turn}
\begin{itemize}
\raggedright
  \item \emph{Say it:} \emph{taḷṇe}
  \item \emph{Say it:} \emph{taḷṇe}, once more
  \item \emph{Say it:} \emph{pāhuṇyāṁsāṭhī ājīne garam pāpaḍ taḷle.} — Grandmother fried hot papads for the guests.
  \item \emph{Recall:} say the Marathi for to chop, then the Marathi for to boil, then say \emph{taḷṇe} again
\end{itemize}

\begin{tcolorbox}[breakable,colback=teal!4,colframe=teal!35!black,title={Before you move on}]
What does \mr{तळणे} mean? (\textquotedblleft{}To fry\textquotedblright{}.) Say it once more.
\end{tcolorbox}
\section[\texorpdfstring{\mr{ढवळणे} (ḍhavaḷṇe)}{ढवळणे (ḍhavaḷṇe)}]{\mr{ढवळणे} (ḍhavaḷṇe) — to stir}

\label{lesson:MR-C527-dhavalne}



\begin{quote}
Before the new one: say the Marathi for to boil, then the Marathi for to fry.
\end{quote}

\subsection*{You'll want to know: \mr{ढवळणे}}
\textbf{\mr{ढवळणे}} — \emph{ḍhavaḷṇe} — \textquotedblleft{}to stir\textquotedblright{}.

\begin{quote}
\mr{आमटी} \mr{उकळताना} \mr{मी} \mr{ती} \mr{चमच्याने} \mr{ढवळली}.
\end{quote}

\begin{quote}
\emph{āmṭī ukaḷtānā mī tī camcyāne ḍhavaḷlī.}
\end{quote}

\begin{quote}
While the dal was boiling I stirred it with a spoon.
\end{quote}

Past, with -\mr{ने} (he …): \mr{त्याने} \mr{आमटी} \mr{ढवळली}, \emph{tyāne āmṭī ḍhavaḷlī}.

\begin{cousinweb}
\textbf{\mr{ढवळणे}} (\emph{ḍhavaḷṇe}) is a Marathi verb for stirring; its origin is not settled.
\end{cousinweb}

\begin{grammarlens}[title={What you've built}]
One more word for describing what you lived and felt.
\end{grammarlens}

\subsection*{Your turn}
\begin{itemize}
\raggedright
  \item \emph{Say it:} \emph{ḍhavaḷṇe}
  \item \emph{Say it:} \emph{ḍhavaḷṇe}, once more
  \item \emph{Say it:} \emph{āmṭī ukaḷtānā mī tī camcyāne ḍhavaḷlī.} — While the dal was boiling I stirred it with a spoon.
  \item \emph{Recall:} say the Marathi for to boil, then the Marathi for to fry, then say \emph{ḍhavaḷṇe} again
\end{itemize}

\begin{tcolorbox}[breakable,colback=teal!4,colframe=teal!35!black,title={Before you move on}]
What does \mr{ढवळणे} mean? (\textquotedblleft{}To stir\textquotedblright{}.) Say it once more.
\end{tcolorbox}
\section[\texorpdfstring{\mr{करपणे} (karapṇe)}{करपणे (karapṇe)}]{\mr{करपणे} (karapṇe) — to get burnt}

\label{lesson:MR-C527-karapne}



\begin{quote}
Before the new one: say the Marathi for to fry, then the Marathi for to stir.
\end{quote}

\subsection*{You'll want to know: \mr{करपणे}}
\textbf{\mr{करपणे}} — \emph{karapṇe} — \textquotedblleft{}to get burnt\textquotedblright{}.

\begin{quote}
\mr{माझं} \mr{लक्ष} \mr{नव्हतं}, \mr{म्हणून} \mr{भात} \mr{करपला}.
\end{quote}

\begin{quote}
\emph{mājhaṁ lakṣa navhataṁ, mhaṇūn bhāt karaplā.}
\end{quote}

\begin{quote}
I was not paying attention, so the rice got burnt.
\end{quote}

Past (he …): \mr{तो} \mr{भात} \mr{थोडा} \mr{करपला}, \emph{to bhāt thoḍā karaplā}.

\mr{करपणे} is for food that burns in the pan.

\begin{cousinweb}
\textbf{\mr{करपणे}} (\emph{karapṇe}) is a Marathi verb for food that scorches in the pan; its origin is not settled.
\end{cousinweb}

\begin{grammarlens}[title={What you've built}]
One more word for describing what you lived and felt.
\end{grammarlens}

\subsection*{Your turn}
\begin{itemize}
\raggedright
  \item \emph{Say it:} \emph{karapṇe}
  \item \emph{Say it:} \emph{karapṇe}, once more
  \item \emph{Say it:} \emph{mājhaṁ lakṣa navhataṁ, mhaṇūn bhāt karaplā.} — I was not paying attention, so the rice got burnt.
  \item \emph{Recall:} say the Marathi for to fry, then the Marathi for to stir, then say \emph{karapṇe} again
\end{itemize}

\begin{tcolorbox}[breakable,colback=teal!4,colframe=teal!35!black,title={Before you move on}]
What does \mr{करपणे} mean? (\textquotedblleft{}To get burnt\textquotedblright{}.) Say it once more.
\end{tcolorbox}
\section[(dialogue)]{Looking back: from Pipes and Leaks to At the Stove}

\label{lesson:MR-R527-review-520-527}



\begin{quote}
Say the Marathi for a faucet and to drip.
\end{quote}

\subsection*{Your turn}
Say these aloud:

\begin{itemize}
\raggedright
  \item a faucet, to drip, a leak, a drain, to get blocked
  \item a fissure, to crack, glue, a saw, to plaster
  \item a living room, a bedroom, a terrace, a floor tile, to hang up
  \item a carpet, a quilt, a mosquito net, a mat, to roll up
  \item rubbish, a grass broom, a floor cloth, to shake off, to scrub
  \item a brass water pot, a clay water pot, to pour, to spill, to rinse
  \item a stove, a kitchen counter, a rolling board, to knead, to grate
  \item to chop, to boil, to fry, to stir, to get burnt
\end{itemize}

\begin{tcolorbox}[breakable,colback=teal!4,colframe=teal!35!black,title={Before you move on}]
A faucet? (\textbf{\mr{नळ}}, \emph{naḷ}.) To drip? (\textbf{\mr{गळणे}}, \emph{gaḷṇe}.) A leak? (\textbf{\mr{गळती}}, \emph{gaḷtī}.) A drain? (\textbf{\mr{गटार}}, \emph{gaṭār}.) To get blocked? (\textbf{\mr{तुंबणे}}, \emph{tumbṇe}.) A fissure? (\textbf{\mr{भेग}}, \emph{bheg}.) To crack? (\textbf{\mr{तडकणे}}, \emph{taḍakṇe}.) Glue? (\textbf{\mr{गोंद}}, \emph{gond}.) A saw? (\textbf{\mr{करवत}}, \emph{karvat}.) To plaster? (\textbf{\mr{लिंपणे}}, \emph{limpṇe}.) A living room? (\textbf{\mr{दिवाणखाना}}, \emph{divāṇkhānā}.) A bedroom? (\textbf{\mr{शयनगृह}}, \emph{śayangṛh}.) A terrace? (\textbf{\mr{गच्ची}}, \emph{gaccī}.) A floor tile? (\textbf{\mr{फरशी}}, \emph{pharśī}.) To hang up? (\textbf{\mr{टांगणे}}, \emph{ṭāṅgṇe}.) A carpet? (\textbf{\mr{गालिचा}}, \emph{gālicā}.) A quilt? (\textbf{\mr{रजई}}, \emph{rajaī}.) A mosquito net? (\textbf{\mr{मच्छरदाणी}}, \emph{macchardāṇī}.) A mat? (\textbf{\mr{चटई}}, \emph{caṭaī}.) To roll up? (\textbf{\mr{गुंडाळणे}}, \emph{guṇḍāḷṇe}.) Rubbish? (\textbf{\mr{कचरा}}, \emph{kacrā}.) A grass broom? (\textbf{\mr{केरसुणी}}, \emph{kersuṇī}.) A floor cloth? (\textbf{\mr{पोतेरे}}, \emph{potere}.) To shake off? (\textbf{\mr{झटकणे}}, \emph{jhaṭakṇe}.) To scrub? (\textbf{\mr{घासणे}}, \emph{ghāsṇe}.) A brass water pot? (\textbf{\mr{कळशी}}, \emph{kaḷśī}.) A clay water pot? (\textbf{\mr{माठ}}, \emph{māṭh}.) To pour? (\textbf{\mr{ओतणे}}, \emph{otṇe}.) To spill? (\textbf{\mr{सांडणे}}, \emph{sāṇḍṇe}.) To rinse? (\textbf{\mr{विसळणे}}, \emph{visaḷṇe}.) A stove? (\textbf{\mr{शेगडी}}, \emph{śegḍī}.) A kitchen counter? (\textbf{\mr{ओटा}}, \emph{oṭā}.) A rolling board? (\textbf{\mr{पोळपाट}}, \emph{poḷpāṭ}.) To knead? (\textbf{\mr{मळणे}}, \emph{maḷṇe}.) To grate? (\textbf{\mr{किसणे}}, \emph{kisṇe}.) To chop? (\textbf{\mr{चिरणे}}, \emph{cirṇe}.) To boil? (\textbf{\mr{उकळणे}}, \emph{ukaḷṇe}.) To fry? (\textbf{\mr{तळणे}}, \emph{taḷṇe}.) To stir? (\textbf{\mr{ढवळणे}}, \emph{ḍhavaḷṇe}.) To get burnt? (\textbf{\mr{करपणे}}, \emph{karapṇe}.)
\end{tcolorbox}
